% ::::::::::::::::::::::::::::::::::::::::::::::::::::::::::::::::::::::::::::::::::
% DETERMINISTIC COLOCATION: ELIMINATING DIGITAL MUDA IN MODERN EXECUTION RUNTIMES
% arXiv submission draft (2026-11-15). Aggressive-track timeline.
% Categories: cs.DC ; cs.SE
% Compile: pdflatex (arxiv-classic). No external fonts, no bibtex: pure LaTeX.
% ::::::::::::::::::::::::::::::::::::::::::::::::::::::::::::::::::::::::::::::::::
\documentclass[10pt,conference]{IEEEtran}
\usepackage[utf8]{inputenc}
\usepackage[T1]{fontenc}
\usepackage{amsmath,booktabs,array}
\title{Deterministic Colocation: Eliminating Digital Muda in Modern Execution Runtimes}
\comments{arXiv submission candidate. Public prior-art reference for the RouteOne Cloud TC charter (pending OASIS TC formation). Repository: READ; benchmark data: benchmark/.}
\author{\IEEEauthorblockN{A. S. Hotta}
\IEEEauthorblockA{Institute for Autonomous Cybernetics \& Applied Monoliths\\
\texttt{hotta@routeone-research.org}}}
\date{15 November 2026}

\begin{document}
\begin{abstract}
Modern enterprise software engineering has conflated operational scale with
computational complexity. Decomposing straightforward domain workflows into
asynchronous, distributed microservice topologies has reintroduced the classic
manifestations of \emph{muda} (waste) identified in the Toyota Production
System---most notably transport latency, queue inventory,
over-processing via repetitive serialization, and defect amplification across
network boundaries. We introduce \emph{Deterministic Colocation
(RouteOne-S)}: an execution paradigm that collapses intermediate network hops
into unified, vertically-integrated pipelines running directly against
sovereign bare-metal OS primitives. Removing distributed consensus layers,
serialization overhead, and ingress virtualization, RouteOne-S achieves a
$94\%$ reduction in CPU idle-cycle waste and collapses p99 tail latencies from
$320\,$ms to $4.2\,$ms. The paradigm demonstrates that runtime resilience is
inversely proportional to topology depth. All data and benchmark harness are
released with this manuscript.
\end{abstract}

% --- body: 8 sections max, IEEE two-column squeeze ---

\section{Introduction: The Pathology of Distributed Overhead}
Modern enterprise configurations exhibit severe mechanical inefficiency: a
single customer request frequently traverses up to twelve internal hops across
reverse proxies, service discovery daemons, API gateways, sidecar containers,
and message queues before touching persistent storage. Every network boundary
crossed introduces (1)~kernel-to-user context switching;
(2)~data duplication across process heaps; (3)~serialization churn---converting
binary states to text, parsing strings back to memory, and repeating the cycle.
We argue structural fragmentation is an anti-pattern and present
\emph{Deterministic Colocation} as a methodology for reclaiming mechanical
sympathy.

\section{Taxonomy of Software Muda}
Adapting the Taiichi Ohno classification to software execution runtimes yields
seven computational failure modes:
\begin{table}[t]
\centering
\begin{tabular}{@{}llp{5.6cm}@{}}
\toprule
\textbf{Classic TPS Muda} & \textbf{Software Equivalent} & \textbf{Manifestation in contemporary stacks} \\
\midrule
Overproduction & Speculative ingestion \& polling & Microservices polling empty queues \\
Waiting & Thread starvation \& I/O stalls & Futures awaiting gRPC/REST responses \\
Transport & Inter-process network hops & Marshalling over loopback/CNI plugins \\
Over-processing & Repetitive SerDe cycles & JSON/Protobuf round-trips $N\times$ per pipeline \\
Inventory & Queue \& buffer saturation & Kafka topics, Redis buffers, socket backlogs \\
Motion & Context switching \& cache thrash & Bloat, saturated dependency trees \\
Defects & Cascading network partitions & Retry storms, misconfigured breakers \\
\bottomrule
\end{tabular}
\caption{Seven waste classes mapped from Toyota Production to software runtime failure modes.}
\label{tab:muda}
\end{table}

\section{The RouteOne-S Method}
RouteOne-S treats the host environment as a single, deterministic assembly
line with three structural constraints: (i)~zero-copy ingress at the native
runtime boundary; (ii)~in-process dispatch rather than RPC; (iii)~single-writer
state transitions (LMAX Disruptor / event-sourcing style). Formally:
\begin{itemize}
\item $\mathit{latency_{total}} = t_{kernel} + t_{compute} + t_{fsync}
\;\ll\; \sum_{i=1..k}(t_{serialize,i}+t_{transit,i}+t_{deserialize,i}) + t_{consensus}$
for $k$ intermediate hops.
\item Allocation per transaction stays bounded and stack-allocated where
possible; garbage-collection stop-the-world pauses degrade accordingly.
\item Disk writes reduce to sequential append I/O; no distributed locks, no
two-phase commit for simple state changes.
\end{itemize}

\section{Empirical Evaluation}
\textbf{Testbed.} Identical bare-metal hardware: AMD EPYC 7763 (64 cores),
256 GiB RAM, NVMe. Mesh baseline: 4 containerised Go microservices behind Envoy
on Kubernetes. RouteOne-S: single process, direct syscalls.

\begin{table}[t]
\centering
\begin{tabular}{@{}lccc@{}}
\toprule
\textbf{Architecture} & \textbf{Throughput} & \textbf{p99 latency} & \textbf{RAM} \\
\midrule
Mesh microservices & 22.4 k req/s & 320 ms & 34.2 GB \\
RouteOne-S monolith & 418 k req/s & 4.2 ms & 1.8 GB \\
\bottomrule
\end{tabular}
\caption{Headline benchmark. Data: \texttt{benchmark/bench-results.csv}.}
\label{tab:bench}
\end{table}

Hardware performance-counter analysis (\texttt{perf stat}) shows the mesh
baseline devoting $\approx$68\% of CPU cycles to socket handling, string
parsing, TLS renegotiation, and memory management; in RouteOne-S, active
compute represents 89\% of cycles.

\section{Related Work and Theological Divergence}
Microservice frameworks (Spring Cloud, Dapr) assume failures are best managed
by dynamic orchestration; actor systems (Akka, Orleans) correctly target
colocation but incur lifecycle overhead when strictly unbounded. RouteOne-S
re-applies the Unix philosophy of small tools connected by \emph{direct
in-process} pipelines, rejecting the modern conflation of overlay IP networks
with \texttt{stdin}/\texttt{stdout}.

\section{Conclusion}
Digital \emph{muda} is rarely a hardware limitation; it is an architectural
choice born of premature distribution. Contemporary x86 and ARM processors
handle high-throughput enterprise workloads within a single coherent execution
context. The highest-performing network call is the one that was never made.

% --------------------
\end{document}
